\subsection{What the Fine-Tuning Actually Did}
\label{sec:retention}

A null result invites the objection that nothing happened at all: that the
adapter never took, or that LoRA at rank 16 is too small to move a model. The
item-level data rules that out. A great deal happened. It simply was not
learning to reason.

Table~\ref{tab:retention} reports answer retention, meaning the share of items
a model answered correctly before fine-tuning that it still answered correctly
afterwards. If 873 traces were teaching mathematical reasoning, retention
should be high and the gains should sit on top of it. That is not the pattern.

\begin{table}[htbp]
\caption{Answer retention. Of the items each model answered correctly
before fine-tuning, how many it still answered correctly afterwards, and
how many items changed outcome in either direction. Notice that no model
retained even half of what it already knew, which is why the small net
changes in Table~\ref{tab:main} understate how much the adapter moved.}
\label{tab:retention}
\centering
\footnotesize
\setlength{\tabcolsep}{4pt}
\begin{tabular}{|l|c|c|c|c|}
\hline
\textbf{Model} & \textbf{Correct} & \textbf{Kept} & \textbf{Reten-} & \textbf{Items} \\
               & \textbf{before}  &               & \textbf{tion}   & \textbf{flipped} \\
\hline
Qwen3-4B & 112 & 56 & 50\% & 99 / 250 \\
Qwen3-1.7B & 42 & 18 & 43\% & 51 / 250 \\
Qwen3-0.6B & 9 & 4 & 44\% & 16 / 250 \\
TigerLLM-1B & 11 & 2 & 18\% & 13 / 250 \\
TituLLM-3B-Instruct & 3 & 0 & 0\% & 6 / 250 \\
TituLLM-1B-Instruct & 1 & 0 & 0\% & 8 / 250 \\
\hline
\end{tabular}
\end{table}

Qwen3-4B answered 112 items correctly before tuning and kept 56 of them. It
also answered 43 items correctly that it had failed before. In total 99 of 250
items changed their outcome, close to 40\% of the evaluation set, while net
accuracy moved by 13 items. The same shape holds down the scale: no model
retained even half of what it already knew. The adapter did not add a layer of
competence on top of the base model. It substantially rewrote which problems
the model could solve, and the replacement set happened to be slightly worse
than the original.

The mechanism is visible in the generated text. Before fine-tuning, Qwen3-4B
closed 225 of 250 completions with the Bengali word for ``answer'' followed by
its result, a convention it adopted on its own, since the evaluation prompt
asks only that the model answer in Bengali and specifies no format. After
fine-tuning it did so on 3 of 250. In its place the model adopted the format of
the training corpus: Ganit's traces are written as a \texttt{<think>} block
followed by an \texttt{<answer>} block, and after tuning every one of the 250
completions opened a \texttt{<think>} block. Median output length fell from
1,046 tokens to 887, the proportion of Bengali characters fell from 87.0\% to
82.0\%, and the number of explicit step markers roughly halved.

This is the Superficial Alignment Hypothesis \parencite{zhou2023lima} observed
from the unfavourable side. In LIMA's setting a small number of examples
exposes capability the pretrained weights already hold, and the result is a
useful model. Here the same small number of examples reshaped format, length,
language mix and stopping behaviour just as comprehensively, and exposed
nothing, because at this scale there was little underneath to expose. Style
transferred completely. Reasoning did not transfer at all.

\subsubsection*{A sensitivity check on the format change}

Because the output format changed, it is worth confirming that the accuracy
drop is not an artifact of answer extraction. Our scorer looks for an
\texttt{<answer>} block first, then a boxed expression, then the Bengali answer
marker, and only if none of those are present does it fall back to taking the
last number in the completion. After fine-tuning, 230 of 250 Qwen3-4B
completions contained a well-formed \texttt{<answer>} block and were therefore
parsed by the strongest branch of the extractor. The remaining 20 exhausted the
generation budget inside the \texttt{<think>} block and were scored by the weak
fallback. Of the 56 items Qwen3-4B lost, 45 produced a well-formed
\texttt{<answer>} that was simply wrong.

To bound the effect, we recomputed accuracy crediting every fallback-scored
item as correct, which is the most generous assumption available. Post-tuning
accuracy under that assumption is 46.8\%, against a measured 39.6\% and a
pre-tuning 44.8\%. The true value therefore lies somewhere between a loss of
5.2 points and a gain of 2.0 points. We cannot sign the effect, which is
consistent with the exact McNemar test returning $p=0.228$, and the substantive
conclusion is unchanged in either direction: 873 curated traces did not produce
a detectable improvement in Bengali mathematical reasoning at this scale.
